\chapter{Estado del Arte}


\section{Algoritmos de Aprendizaje por Refuerzo para Control Reactivo}
El desarrollo del DRL aplicado al guiado de robots autónomos se divide primordialmente entre algoritmos de optimización de gradiente de política directos \textit{on-policy} y algoritmos basados en valor de acción de tipo crítico-actor \textit{off-policy} \citep{perez2022navigation}. Entre los enfoques \textit{on-policy}, el algoritmo de Optimización de Política Próxima (PPO, por sus siglas en inglés) destaca por su extraordinaria estabilidad numérica y robustez ante el ajuste de hiperparámetros \citep{RL_Algorithms_Brief.pdf}. PPO previene las actualizaciones de política excesivamente destructivas mediante una función de pérdida sustituta con restricciones por recorte del ratio de probabilidad de la acción, limitando la divergencia de Kullback-Leibler entre la política vieja y la nueva \citep{2507.05251v1.pdf}:
\begin{equation}
    L^{CLIP}(\theta) = \hat{\mathbb{E}}_t \left[ \min(r_t(\theta)\hat{A}_t, \text{clip}(r_t(\theta), 1-\epsilon, 1+\epsilon)\hat{A}_t) \right]
\end{equation}
No obstante, dado que PPO es \textit{on-policy}, padece de una baja eficiencia en el uso de muestras, requiriendo millones de interacciones directas con el entorno de simulación para converger hacia una política robusta \citep{state_art_RL.pdf}.

Para solventar la ineficiencia de muestreo, se han extendido aproximaciones \textit{off-policy} de crítico-actor continuas como DDPG y su variante de críticos retrasados gemelos TD3 \citep{sivashangaran2023deep}. Adicionalmente, el algoritmo Soft Actor-Critic (SAC) destaca como un método altamente eficiente basado en la maximización de la entropía máxima, el cual optimiza de manera conjunta el retorno esperado y el grado de aleatoriedad de la política \citep{DRL_AGV_Exploration.pdf}. Esta adición fomenta una exploración muy agresiva y previene que el actor converja prematuramente a mínimos locales estáticos o cuellos de botella subóptimos en la navegación \citep{applsci-12-03153-v2.pdf}:
\begin{equation}
    J(\pi) = \sum_{t=0}^{T} \mathbb{E}_{(s_t, a_t) \sim \rho_\pi} \left[ R(s_t, a_t) + \alpha \mathcal{H}(\pi(\cdot | s_t)) \right]
\end{equation}
En entornos logísticos de rejilla donde las acciones son intrínsecamente discretas (avanzar, girar), las redes Q profundas (DQN) de tipo \textit{off-policy} y sus extensiones de doble estimación (Double DQN) resultan óptimas por su simplicidad de convergencia teórica y por permitir el reuso masivo de muestras previas del buffer de replay \citep{hu2020deep}.

\section{Percepción Sensorial Multimodal y Observabilidad Parcial}
La navegación reactiva sin mapas se enmarca típicamente bajo la formulación matemática de un Proceso de Decisión de Markov Parcialmente Observable (POMDP), dado que el AGV solo dispone de lecturas de sensores de rango locales y carece de la posición global de los obstáculos en un plano cartesiano unificado \citep{md2024novel}. El sensado local con LiDARs discretos o continuos presenta problemas críticos de oclusión geométrica (puntos ciegos detrás de paredes) y ruidos sensoriales de rebote estocástico \citep{sivashangaran2023autovrl}. Para atenuar la dimensionalidad de las observaciones crudas de LiDAR (que pueden llegar a cientos de rayos densos), la literatura propone técnicas de discretización modular o compresión mediante autoencoders geométricos robustos (GDAE) \citep{article_1743243135.pdf}.

Por otra parte, la fusión de información sensorial visual profunda con datos de rango LiDAR se consolida mediante el uso de redes convolucionales preentrenadas (como MobileNetV3) en esquemas autoasociativos estáticos para codificar el entorno visual en un vector de baja dimensionalidad espacial \citep{2507.05251v1.pdf, applsci-12-03153-v2.pdf}. No obstante, la pérdida de observabilidad temporal en el POMDP requiere el uso de memorias recurrentes profundas, tales como celdas LSTM o arquitecturas autorregresivas basadas en Transformers de decisión (Decision Transformers) \citep{ge2024deep, fnbot-18-1338189.pdf}. Estos componentes de memoria temporal permiten procesar una ventana histórica de observaciones y acciones previas $h_t = (o_1, a_1, \dots, o_t)$, extrayendo el ``estado de creencia'' estocástico implícito (\textit{belief state}) indispensable para desambiguar estados geométricamente idénticos y evadir de forma activa obstáculos dinámicos u oclusiones severas \citep{md2024novel, state_art_RL.pdf}.

\section{Aprendizaje por Currículo y Navegación Segura}
El entrenamiento de un AGV map-less desde cero bajo un esquema tradicional de DRL se enfrenta al cuello de botella de la ``recompensa dispersa'' (\textit{sparse reward}) \citep{ye2023toward}. Si el espacio de estados es grande y la meta es de dimensiones reducidas en comparación con la arena de pruebas, la probabilidad estocástica de que el agente alcance la meta mediante exploración aleatoria pura al inicio del entrenamiento es prácticamente nula, lo que provoca estancamientos en el gradiente de optimización \citep{state_art_RL.pdf}. Para mitigar esto, se ha introducido el Aprendizaje por Currículo Automático (ACL, por sus siglas en inglés) \citep{applsci-12-03153-v2.pdf}.

La metodología propuesta por el algoritmo \textit{NavACL} y su versión extendida con retroalimentación del valor de acción de los críticos, \textit{NavACL-Q}, proponen dinámicamente escenarios de entrenamiento de dificultad adaptativa progresiva \citep{applsci-12-03153-v2.pdf}. NavACL-Q evalúa métricas geométricas de los mapas logísticos (distancia geodésica, holgura del agente y del objetivo ante obstáculos, y el ángulo de rumbo inicial) combinándolas con la estimación de valor Q predicha por la red crítica del agente, para proponer poses iniciales que se sitúen en la frontera del conocimiento actual de la política \citep{xue2022using}. A medida que el AGV incrementa su tasa de éxito local, el currículo de forma automática expande la distancia inicial de navegación y la complejidad de los laberintos. Adicionalmente, el Aprendizaje por Refuerzo Seguro (Safe RL) incorpora la teoría de Optimización de Política con Restricciones (CPO, por sus siglas en inglés), donde la función de costo modela explícitamente el riesgo de colisión catastrófica, asegurando que la política explore el espacio con garantías de seguridad física bajo restricciones de barreras de control \citep{gao2025hierarchical}.
